\section{La disolución de las fronteras: Browser, Desktop, Mobile, GUI, CLI, API y Coding}

Su Yu (苏煜) insiste en que, en sus inicios, el campo de los Agent distinguía entre browser use, desktop use, mobile use, GUI, text-based representation, CLI, API, coding, etc., pero que estas divisiones son provisionales. Lo que finalmente se busca es un universal digital agent. Todas estas interfaces son medios, no el fin último.

\begin{figure}[H]
\centering
\includegraphics[width=0.92\textwidth]{boundary-convergence.png}
\caption{La disolución de las fronteras: múltiples puntos de entrada a entornos digitales están convergiendo hacia un universal digital agent.}
\end{figure}

\begin{knowledgebox}{Lectura de la figura: por qué el coding es el núcleo de la disolución de las fronteras}
En la figura, todos los puntos de entrada se conectan con el universal digital agent. Browser, desktop, mobile y GUI son la interacción de superficie; CLI, API y coding son interfaces de control más formalizadas. Su Yu considera que el coding es el fabric del digital world, porque muchas interfaces pueden, en última instancia, expresarse, renderizarse o manipularse mediante código.
\end{knowledgebox}

\subsection{La GUI no desaparecerá, y la CLI tampoco lo dominará todo}

La entrevista no llega al extremo de afirmar que «la CLI reemplazará por completo a la GUI». Hay dos razones. Primera: en el mundo real existe una gran cantidad de legacy system, por ejemplo en bancos, software empresarial y sistemas antiguos, que no se reescribirán con rapidez. Segunda: aunque la CLI sea el óptimo global para los Agent, eso no significa que sea la mejor opción en todos los escenarios locales. Muchas GUI existentes ya son good enough para las personas y las organizaciones.

\begin{warningbox}{El óptimo local es distinto del óptimo global}
Una interfaz técnicamente más adecuada para los Agent no necesariamente sustituirá de inmediato a las interfaces existentes. Los costos de migración de las empresas, los hábitos de los usuarios, la regulación, los sistemas heredados y las cuentas económicas determinan la trayectoria real. El juicio técnico sobre los Agent debe considerarse junto con el entorno de despliegue.
\end{warningbox}

\subsection{Un Programming Language también es un Language}

El conductor propone la metáfora de que «el lenguaje natural es el andamiaje de los humanos y el coding es el andamiaje de las máquinas». Su Yu añade que language nunca se ha limitado al lenguaje natural: los lenguajes de programación, los diagramas y los gestos son language en sentido amplio. Un programming language es un formal language, más preciso y más ejecutable, y por ello resulta especialmente importante cuando un Agent opera en el digital world.

\subsection{Resumen de la sección}

El desenlace de los Agent no es la victoria individual del browser agent, el desktop agent o el coding agent, sino la convergencia de múltiples puntos de entrada en el nivel de la tarea. La importancia del coding radica en que unifica muchas interfaces como objetos que pueden expresarse, combinarse y ejecutarse.
